% !TEX program = xelatex
\documentclass[12pt,a4paper]{article}

% ============================================================================
% Packages & Language Setup (XeLaTeX + Thai)
% ============================================================================
\usepackage{geometry}
\geometry{
    a4paper,
    top=2.5cm,
    bottom=2.5cm,
    left=2.5cm,
    right=2.5cm,
    headheight=16pt
}

\usepackage{fontspec}
\usepackage{polyglossia}
\setdefaultlanguage{thai}
\setotherlanguage{english}

% ฟอนต์ภาษาไทยมาตรฐานรายงาน (โหลดจากโฟลเดอร์ docs/fonts/ เพื่อความสอดคล้องทุกเครื่อง)
\setmainfont[
    Path = fonts/,
    Extension = .ttf,
    Script = Thai,
    Scale = 1.35,
    UprightFont = THSarabunNew,
    BoldFont = THSarabunNew-Bold,
    ItalicFont = THSarabunNew-Italic,
    BoldItalicFont = THSarabunNew-BoldItalic
]{THSarabunNew}

\newfontfamily\thaifont[
    Path = fonts/,
    Extension = .ttf,
    Script = Thai,
    Scale = 1.35,
    UprightFont = THSarabunNew,
    BoldFont = THSarabunNew-Bold,
    ItalicFont = THSarabunNew-Italic,
    BoldItalicFont = THSarabunNew-BoldItalic
]{THSarabunNew}

\IfFontExistsTF{Times New Roman}{%
    \newfontfamily\englishfont[Scale=1.0]{Times New Roman}%
}{%
    \newfontfamily\englishfont[Scale=1.0]{Latin Modern Roman}%
}
\IfFontExistsTF{Courier New}{%
    \newfontfamily\thaifonttt[Script=Thai, Scale=1.0]{Courier New}%
    \setmonofont[Scale=0.95]{Courier New}%
}{%
    \newfontfamily\thaifonttt[Scale=1.0]{Latin Modern Mono}%
    \setmonofont[Scale=0.95]{Latin Modern Mono}%
}

% การตัดคำภาษาไทยอัตโนมัติใน XeLaTeX
\XeTeXlinebreaklocale "th"
\XeTeXlinebreakskip = 0pt plus 1pt

\usepackage{graphicx}
% พาธสำหรับรูปภาพและสไปรต์ที่ scale up แล้วใน docs/images/ และ docs/images/sprite/
\graphicspath{{images/}{images/sprite/}{../../}{./}}

\usepackage{amsmath}
\usepackage{booktabs}
\usepackage{array}
\usepackage{tabularx}
\usepackage{longtable}
\usepackage{float}
\usepackage{caption}
\usepackage{enumitem}
\usepackage{fancyhdr}
\usepackage{titlesec}
\usepackage{xcolor}
\usepackage{listings}
\usepackage{hyperref}

\newcolumntype{L}{>{\raggedright\arraybackslash}X}
\captionsetup{justification=centering}

\lstset{
    basicstyle=\ttfamily\small,
    breaklines=true,
    frame=single,
    rulecolor=\color{gray!40},
    backgroundcolor=\color{gray!5},
    numbers=none,
    tabsize=4,
    showstringspaces=false,
    columns=flexible,
    xleftmargin=0.5em,
    xrightmargin=0.5em
}

\hypersetup{
    colorlinks=true,
    linkcolor=black,
    urlcolor=blue,
    pdftitle={รายงานผลการดำเนินการโครงงาน การพัฒนาเกม SLIME MAZE ด้วย FPGA},
    pdfauthor={คณะผู้จัดทำโครงงาน CPE222}
}

% ============================================================================
% Heading & TOC Numbering Style (คล้าย REF_TH_CPE222_Final_Documentation)
%   - หัวข้อหลัก (Section): ไม่มีตัวเลขนำหน้า แต่แสดงในสารบัญ
%   - หัวข้อย่อยระดับ 1 (Subsection): 1), 2), 3) ...
%   - หัวข้อย่อยระดับ 2 (Subsubsection): 1.1), 1.2), ...
%   - หัวข้อย่อยระดับ 3 (Paragraph): 1.1.1), 1.1.2), ...
% ============================================================================
\setcounter{secnumdepth}{4}
\setcounter{tocdepth}{4}

\renewcommand{\thesubsection}{\arabic{subsection})}
\renewcommand{\thesubsubsection}{\arabic{subsection}.\arabic{subsubsection})}
\renewcommand{\theparagraph}{\arabic{subsection}.\arabic{subsubsection}.\arabic{paragraph})}

% ปรับระยะเยื้องในสารบัญให้พอดี ไม่ล้นขอบขวา
\makeatletter
\renewcommand*\l@section{\@dottedtocline{1}{0em}{0em}}
\renewcommand*\l@subsection{\@dottedtocline{2}{1.2em}{2.0em}}
\renewcommand*\l@subsubsection{\@dottedtocline{3}{2.8em}{2.8em}}
\renewcommand*\l@paragraph{\@dottedtocline{4}{4.6em}{3.6em}}
\makeatother

% คำสั่งสำหรับหัวข้อหลักที่ไม่มีเลขนำหน้า แต่ขึ้นสารบัญและรีเซ็ตเลขหัวข้อย่อย
\newcommand{\mainsection}[1]{%
    \phantomsection
    \stepcounter{section}%
    \section*{#1}%
    \addcontentsline{toc}{section}{#1}%
}

% ข้อความรอเติมเนื้อหา (ลบออกได้เมื่อเริ่มเขียนเนื้อหาแต่ละหัวข้อ)
\newcommand{\todo}[1]{\noindent\textcolor{gray}{\small\textit{[รอเพิ่มเนื้อหา: #1]}}\par\vspace{0.3em}}

% กล่องแสดงตำแหน่งสำหรับใส่ภาพ (ใช้เมื่อยังไม่ได้วางไฟล์ภาพจริง)
\newcommand{\imgplaceholder}[3]{%
    \begin{center}
        \fcolorbox{gray!60}{gray!6}{%
            \parbox[c][#1][c]{0.92\linewidth}{%
                \centering\small\color{gray!85!black}%
                \textbf{[พื้นที่สำหรับใส่ภาพ: \texttt{\detokenize{#2}}]}\\[0.25em]
                #3%
            }%
        }%
    \end{center}%
}

% แทรกภาพอัตโนมัติจากพาธมาตรฐาน หากยังไม่มีไฟล์ภาพจะแสดงกล่องคำอธิบายแทน (ไม่ทำให้คอมไพล์พัง)
\newcommand{\smartimage}[4]{%
    \IfFileExists{images/#2}{%
        \includegraphics[width=#1]{images/#2}%
    }{%
        \IfFileExists{images/sprite/#2}{%
            \includegraphics[width=#1]{images/sprite/#2}%
        }{%
            \IfFileExists{../../#2}{%
                \includegraphics[width=#1]{../../#2}%
            }{%
                \IfFileExists{#2}{%
                    \includegraphics[width=#1]{#2}%
                }{%
                    \imgplaceholder{#3}{#2}{#4}%
                }%
            }%
        }%
    }%
}

% คำบรรยายรูปภาพภายใน longtable (หลีกเลี่ยง \caption ของ longtable ที่ชนกับ \captionof)
\newcommand{\tablefigcap}[1]{%
    \refstepcounter{figure}%
    \addcontentsline{lof}{figure}{\protect\numberline{\thefigure}{\ignorespaces #1}}%
    \par\vspace{0.3em}{\small\textbf{รูปที่ \thefigure}\ #1}\par%
}

\titleformat{\section}{\Large\bfseries}{}{0pt}{}
\titleformat{\subsection}{\large\bfseries}{\thesubsection}{0.6em}{}
\titleformat{\subsubsection}{\normalsize\bfseries}{\thesubsubsection}{0.6em}{}
\titleformat{\paragraph}[block]{\normalsize\bfseries}{\theparagraph}{0.6em}{}

\titlespacing*{\section}{0pt}{1.2em}{0.4em}
\titlespacing*{\subsection}{0pt}{0.8em}{0.2em}
\titlespacing*{\subsubsection}{0.6cm}{0.6em}{0.2em}
\titlespacing*{\paragraph}{1.2cm}{0.5em}{0.2em}

% ชื่อสารบัญภาษาไทย (ตั้งค่าใน \captionsthai ของ polyglossia)
\gappto\captionsthai{%
    \renewcommand{\contentsname}{สารบัญ}%
    \renewcommand{\listfigurename}{สารบัญรูปภาพ}%
    \renewcommand{\listtablename}{สารบัญตาราง}%
    \renewcommand{\figurename}{รูปที่}%
    \renewcommand{\tablename}{ตาราง}%
}

% เลขหน้าอักษรไทย (ก, ข, ค, ...) สำหรับส่วนหน้า
\makeatletter
\def\@thaialph#1{%
    \ifcase#1\or ก\or ข\or ค\or ง\or จ\or ฉ\or ช\or ซ\or ฌ\or ญ%
    \or ฎ\or ฏ\or ฐ\or ฑ\or ฒ\or ณ\or ด\or ต\or ถ\or ท%
    \or ธ\or น\or บ\or ป\or ผ\or ฝ\or พ\or ฟ\or ภ\or ม%
    \else\@ctrerr\fi}
\def\thaialph#1{\expandafter\@thaialph\csname c@#1\endcsname}
\makeatother

% Header / Footer (เลขหน้ามุมขวาบน ตามรูปแบบรายงานอ้างอิง)
\pagestyle{fancy}
\fancyhf{}
\fancyhead[R]{\thepage}
\renewcommand{\headrulewidth}{0pt}
\renewcommand{\footrulewidth}{0pt}

\setlength{\parindent}{1.2cm}
\setlength{\parskip}{0.3em}

% ============================================================================
% Document Start
% ============================================================================
\begin{document}

% หน้าปก
% ----------------------------------------------------------------------------
% หน้าปก (Cover Page)
% ----------------------------------------------------------------------------
\begin{titlepage}
    \centering
    \vspace*{0.5cm}

    % โลโก้มหาวิทยาลัย (ไฟล์ docs/images/KMUTT_Logo.png)
    \IfFileExists{images/KMUTT_Logo.png}{%
        \includegraphics[width=3.2cm]{KMUTT_Logo.png}\\[0.8cm]
    }{\vspace*{2.5cm}}

    {\Large\bfseries รายงานผลการดำเนินการโครงงาน (Final Project)\\[0.3cm]
    การพัฒนาเกม \textenglish{SLIME MAZE} ด้วย \textenglish{FPGA}\par}

    \vspace{1.5cm}

    {\large\bfseries คณะผู้จัดทำ\par}
    \vspace{0.3cm}
    \begin{tabular}{@{}ll@{}}
        1. นาย กฤษติธี อัศวดิษฐเลิศ & \textenglish{68070501002} \\
        2. นาย วรเมธ กิตติโชติรัตน์ & \textenglish{68070501045} \\
        3. นาย อรรถเชษฐ์ อนันต์คูศรี & \textenglish{68070501059} \\
        4. นาย ปริญญา เส็มเส็น & \textenglish{68070501073} \\
        5. นาย พัชรพล ธรรมมล & \textenglish{68070501077} \\
    \end{tabular}

    \vspace{1.5cm}

    {\large\bfseries เสนอ\par}
    \vspace{0.2cm}
    {\large ดร.ทวีชัย นันทวิสุทธิวงศ์\par}

    \vfill

    {\large
    ภาคเรียนที่ 1 ปีการศึกษา 2569 \\
    รายวิชา \textenglish{CPE222} อิเล็กทรอนิกส์ดิจิทัลและการออกแบบวงจรเชิงตรรกะ \\
    \textenglish{(Digital Electronics and Logic Design)} \\
    ภาควิชาวิศวกรรมคอมพิวเตอร์ คณะวิศวกรรมศาสตร์ \\
    มหาวิทยาลัยเทคโนโลยีพระจอมเกล้าธนบุรี
    \par}
\end{titlepage}


% ----------------------------------------------------------------------------
% ส่วนหน้า: สารบัญ, สารบัญรูปภาพ, สารบัญตาราง (เลขหน้า ก, ข, ...)
% ----------------------------------------------------------------------------
\pagenumbering{thaialph}
\setcounter{page}{1}

\phantomsection
\addcontentsline{toc}{section}{สารบัญ}
\tableofcontents
\clearpage

\phantomsection
\addcontentsline{toc}{section}{สารบัญรูปภาพ}
\listoffigures
\vspace{1.5cm}

\phantomsection
\addcontentsline{toc}{section}{สารบัญตาราง}
\listoftables
\clearpage

% ----------------------------------------------------------------------------
% ส่วนเนื้อหา (เลขหน้า 1, 2, 3, ...)
% ----------------------------------------------------------------------------
\pagenumbering{arabic}
\setcounter{page}{1}

% ============================================================================
\mainsection{บทนำ (Introduction)}
% ============================================================================

โครงงาน \textenglish{Slime Maze} คือโครงงานการประยุกต์ใช้ความรู้ในรายวิชา \textenglish{CPE222} อิเล็กทรอนิกส์ดิจิทัลและการออกแบบวงจรเชิงตรรกะ (Digital Electronics and Logic Design) มาออกแบบและพัฒนาชิ้นงานจริงบนฮาร์ดแวร์ FPGA โดยคณะผู้จัดทำมีความสนใจในการพัฒนาระบบเกมอาร์เคดที่มีแนวคิดต้นแบบมาจากเกมแนวเขาวงกตไล่ล่า (Maze-Chase Arcade Game) อย่าง \textenglish{Pac-Man} ผสานกับระบบการวิวัฒนาการของตัวละครสไลม์ 3 ขั้น (Slime Evolution) ที่พัฒนาด้วยภาษา Verilog HDL และแสดงผลผ่านพอร์ต VGA ความละเอียด $640\times480$ พิกเซล บนบอร์ด Digilent Nexys A7-100T โดยมีรายละเอียดพื้นฐานที่เกี่ยวข้องดังนี้

\begin{figure}[H]
    \centering
    \smartimage{0.82\textwidth}{banner.png}{4.0cm}{ภาพแบนเนอร์แนะนำเกม Slime Maze (images/banner.png)}
    \caption{ภาพแสดงผลหน้าปกและแนวคิดของเกม Slime Maze}
\end{figure}

\subsection{แนวคิดและแรงบันดาลใจของเกมเขาวงกต (Arcade Maze Game)}

เกมแนวเขาวงกตไล่ล่า (Maze-Chase Game) เป็นเกมอาร์เคดคลาสสิกที่มีต้นแบบสำคัญมาจากเกม Pac-Man โดยผู้เล่นจะต้องบังคับตัวละครเคลื่อนที่ไปตามทางเดินในเขาวงกตแบบ 2 มิติ เพื่อเก็บเม็ดคะแนน (Pellets) ให้ครบตามเงื่อนไข พร้อมกับหลบหลีกศัตรู (Ghosts) ที่คอยไล่ล่าจากทิศทางต่างๆ แต่เมื่อเก็บไอเทมพิเศษได้ บทบาทจะพลิกกลับจาก “ผู้ถูกล่า” เป็น “ผู้ล่า” ทำให้สามารถไล่กินศัตรูเพื่อทำคะแนนโบนัสได้ชั่วคราว

โครงงาน Slime Maze ได้นำแนวคิดดังกล่าวมายกระดับด้วยระบบวิวัฒนาการตัวละคร 3 ร่าง (Baby Slime $\rightarrow$ Big Slime $\rightarrow$ King Slime) โดยผู้เล่นต้องเก็บเม็ดพลังงานในเขาวงกตขนาด $20\times15$ ช่อง ให้ครบ 30 เม็ดเพื่อเติบโตจากร่างที่ 1 เป็นร่างที่ 2 จากนั้นมงกุฎ (Crown) จะปรากฏขึ้นตามมุมของแผนที่ เมื่อเก็บมงกุฎได้ สไลม์จะเข้าสู่ “โหมดราชา” (King Mode / ร่างที่ 3) เป็นเวลา 10 วินาที ซึ่งจะเดินเร็วขึ้น 2 เท่า และสามารถไล่กินผีทั้ง 4 ตัวเพื่อเอาชนะเกมได้ หากพลาดถูกผีจับตัวครบ 3 ครั้ง เกมจะจบลงทันที

% [INSERT PICTURE: pacman_ref.png — ภาพหน้าจอเกม Pac-Man ต้นแบบเพื่อเปรียบเทียบแรงบันดาลใจ]
\begin{figure}[H]
    \centering
    \smartimage{0.68\textwidth}{pacman_ref.png}{4.2cm}{ภาพหน้าจอ (Screenshot) ของเกม Pac-Man ต้นแบบ แสดงแผนที่เขาวงกต เม็ดคะแนน ตัวละครผู้เล่น และผี 4 ตัว เพื่อใช้อ้างอิงแรงบันดาลใจในการออกแบบ (เพิ่ม URL ที่มาในบรรณานุกรม)}
    \caption{ตัวอย่างภาพแสดงผลเกมต้นแบบแนวเขาวงกตไล่ล่า (Pac-Man)}
\end{figure}

\subsection{Field Programmable Gate Array (FPGA)}

FPGA เป็นอุปกรณ์วงจรรวม (Integrated Circuit) ที่ผู้ใช้สามารถเขียนโปรแกรมเพื่อกำหนดโครงสร้างวงจรดิจิทัลภายในได้ตามต้องการ ภายในชิป FPGA ประกอบด้วยหน่วยตรรกะพื้นฐานจำนวนมาก เช่น Look-Up Tables (LUTs), Flip-Flops, Block RAM (BRAM) และตัวจัดการสัญญาณนาฬิกา (Clock Management Tiles / MMCM) ที่สามารถเชื่อมต่อกันเป็นวงจรดิจิทัลเชิงตรรกะและวงจรลำดับได้อย่างหลากหลาย เช่น วงจรนับ (Counter), วงจรบวก, เครื่องสถานะจำกัด (Finite State Machine: FSM), หน่วยความจำแบบ ROM และวงจรสร้างสัญญาณภาพ (Video Controller)

จุดเด่นสำคัญของ FPGA คือความสามารถในการประมวลผลแบบคู่ขนาน (Parallel Processing) ในระดับฮาร์ดแวร์โดยไม่ต้องพึ่งพาระบบปฏิบัติการหรือซีพียูแบบซอฟต์แวร์ การพัฒนาวงจรบน FPGA ใช้ภาษาคำอธิบายฮาร์ดแวร์ (Hardware Description Language: HDL) เช่น Verilog หรือ VHDL จากนั้นจึงนำโค้ดผ่านกระบวนการ Synthesis และ Implementation เพื่อสร้างไฟล์กำหนดค่าที่เรียกว่า Bitstream สำหรับโปรแกรมลงบนชิป FPGA ต่อไป

\subsection{Digilent Nexys A7-100T Artix-7 FPGA Trainer Board}

บอร์ดพัฒนา Digilent Nexys A7-100T เป็นบอร์ดฝึกปฏิบัติและพัฒนาวงจรดิจิทัลที่มาพร้อมกับชิป FPGA ตระกูล Xilinx Artix-7 รหัส \texttt{xc7a100tcsg324-1} ซึ่งเป็นรุ่นที่มีความจุทางตรรกะสูงเพียงพอต่อการออกแบบระบบเกมและกราฟิกแบบเต็มรูปแบบ บนตัวบอร์ดติดตั้งอุปกรณ์ต่อพ่วงพื้นฐานมาครบถ้วนพร้อมใช้งานทันทีโดยไม่ต้องต่อวงจรภายนอกเพิ่มเติม ได้แก่

\begin{itemize}[leftmargin=1.5cm]
    \item สัญญาณนาฬิกาหลักบนบอร์ด (On-board Oscillator) ความถี่ 100 MHz (ขา E3)
    \item ปุ่มกดแบบ Pushbutton 5 ทิศทาง (\texttt{BTNU}, \texttt{BTND}, \texttt{BTNL}, \texttt{BTNR}, \texttt{BTNC}) และปุ่มรีเซ็ต \texttt{CPU\_RESETN} (Active-Low)
    \item หลอดไฟ LED แสดงสถานะจำนวน 16 ดวง (\texttt{LED[15:0]})
    \item จอแสดงผลตัวเลข 7-Segment แบบ Common Anode จำนวน 8 หลัก (\texttt{SEG[6:0]}, \texttt{DP}, \texttt{AN[7:0]})
    \item พอร์ตเชื่อมต่อสัญญาณภาพ VGA แบบ 12-bit Color (4-bit Red, 4-bit Green, 4-bit Blue พร้อมสัญญาณ \texttt{VGA\_HS} และ \texttt{VGA\_VS}) ผ่านวงจร R-2R Ladder DAC บนบอร์ด
    \item พอร์ต Micro-USB สำหรับจ่ายไฟเลี้ยงและโปรแกรม Bitstream ผ่านวงจร USB-JTAG
\end{itemize}

% ============================================================================
\mainsection{วัตถุประสงค์}
% ============================================================================
\todo{ระบุวัตถุประสงค์ของโครงงาน เช่น การประยุกต์ใช้วิชา CPE222, การสร้างสัญญาณ VGA, การออกแบบเกมบน FPGA}

% ============================================================================
\mainsection{อุปกรณ์ที่ใช้ (Equipment)}
% ============================================================================

\subsection{ซอฟต์แวร์ (Software)}
\todo{ตารางซอฟต์แวร์ เช่น Vivado Design Suite, Verilator + SDL2, Python 3 (Pillow), โปรแกรมวาดสไปรต์}

\subsection{ฮาร์ดแวร์ (Hardware)}
\todo{ตารางฮาร์ดแวร์ เช่น Digilent Nexys A7-100T, สาย VGA, จอภาพรองรับ VGA 640x480 @ 60 Hz}

% ============================================================================
\mainsection{ส่วนประกอบวงจร}
% ============================================================================
\todo{แผนภาพบล็อกส่วนประกอบของระบบโดยรวม}

\subsection{Input Switch (ปุ่มกดควบคุมบนบอร์ด)}
\todo{ปุ่ม BTNU, BTND, BTNL, BTNR, BTNC และ CPU\_RESETN}

\subsection{Field Programmable Gate Array (FPGA)}
\todo{การทำงานของชิป Artix-7 บนบอร์ด Nexys A7-100T}

\subsection{Monitor Display ที่รองรับ VGA และส่วนแสดงผลบนบอร์ด}
\todo{หลักการแสดงผล VGA (12-bit RGB, H-Sync, V-Sync), 7-Segment 8 หลัก และ LEDs}

% ============================================================================
\clearpage
\mainsection{หลักการทำงาน}
% ============================================================================

\subsection{ระบบรูปแบบการเล่นเกม (Game Mechanics)}

\subsubsection{การควบคุมตัวละคร (Character Control)}
\paragraph{การเคลื่อนที่และการกดเลี้ยวล่วงหน้า (Movement \& Pre-turn Input)}
\todo{อธิบายการเดินตามตาราง 20x15 ช่องและการกดปุ่มเลี้ยวล่วงหน้าก่อนถึงทางแยก}

\paragraph{การเริ่มเกม การหยุดชั่วคราว และการรีเซ็ต (Start, Pause \& Reset)}
\todo{อธิบายหน้าที่ของปุ่ม BTNC และ CPU\_RESETN}

\subsubsection{ระบบการวิวัฒนาการของสไลม์ (Evolution \& King Mode)}
\paragraph{ร่างที่ 1 (Form 1: Small Slime)}
\todo{ร่างเริ่มต้น เก็บเม็ด (+10 คะแนน) และหลบผี 4 ตัว (ภาพสไปรต์จาก images/sprite/slime1\_*.png)}

\paragraph{ร่างที่ 2 (Form 2: Big Slime)}
\todo{เก็บเม็ดครบ 30 เม็ด (N\_EVOLVE) มงกุฎจะปรากฏขึ้นในเขาวงกต (ภาพสไปรต์จาก images/sprite/slime2\_*.png)}

\paragraph{ร่างที่ 3 (Form 3: King Mode)}
\todo{เก็บมงกุฎในร่าง 2 เดินเร็วขึ้น 2 เท่า กินผีได้ (+200 คะแนน/ตัว) นาน 10 วินาที (ภาพสไปรต์จาก images/sprite/slime3\_*.png)}

\subsubsection{ระบบศัตรูและไอเทมภายในเขาวงกต (Enemies \& Items)}
\paragraph{พฤติกรรมของผีทั้ง 4 ตัว (Ghost AI)}
\todo{อธิบายผี 4 ตัว: Ghost 0 (Chase), Ghost 1 \& 3 (Chase/Random), Ghost 2 (Random) และโหมด Fright (images/sprite/ghost\_*.png)}

\paragraph{เม็ดคะแนน (Pellets) และมงกุฎ (Crown)}
\todo{การวางเม็ดคะแนนและการสุ่มเกิดใหม่ของมงกุฎ (CROWN\_DELAY = 180 เฟรม) (images/sprite/ball.png, king0..2.png)}

\subsubsection{ระบบคะแนน พลังชีวิต และเงื่อนไขการแพ้ชนะ (Scoring, Lives \& Win/Loss)}
\todo{คะแนนสูงสุด 9990, พลังชีวิต 3 ดวง (life\_alive.png / life\_died.png), เงื่อนไข WIN และ GAME OVER}

\subsection{ระบบภาพรวมการทำงาน (System Overview)}

\subsubsection{สถาปัตยกรรมระดับบนสุด (\texttt{slime\_top.v})}
\todo{แผนภาพการเชื่อมต่อระหว่าง clk\_gen, vga\_sync, slime\_logic, slime\_draw, seg7x8 และ LEDs}

\subsubsection{สถานะการทำงานของเกม (Game FSM ใน \texttt{slime\_logic.v})}
\paragraph{Title State (\texttt{S\_TITLE})}
\todo{หน้าเริ่มต้น แสดงข้อความ SLIME MAZE และ PRESS BTNC}

\paragraph{Play State (\texttt{S\_PLAY})}
\todo{สถานะการเล่นเกมหลัก การเดิน การเก็บเม็ด และการอัปเดต AI ผี}

\paragraph{Pause State (\texttt{S\_PAUSE})}
\todo{สถานะหยุดเกมชั่วคราวเมื่อกด BTNC ระหว่างเล่น}

\paragraph{Dead State (\texttt{S\_DEAD})}
\todo{สถานะเมื่อสไลม์ถูกผีจับได้ (ลดชีวิต 1 ดวง และรีเซ็ตตำแหน่งเมื่อยังมีชีวิตเหลือ)}

\paragraph{Win State (\texttt{S\_WIN})}
\todo{สถานะชนะเกมเมื่อกินผีครบทั้ง 4 ตัว}

\paragraph{Game Over State (\texttt{S\_OVER})}
\todo{สถานะจบเกมเมื่อพลังชีวิตเหลือ 0}

\subsection{ระบบฮาร์ดแวร์และสัญญาณควบคุม (Hardware System)}

\subsubsection{สัญญาณนาฬิกาและวงจรรีเซ็ต (\texttt{clk\_gen.v} และ Reset Sync)}
\todo{MMCM แปลง 100 MHz -> 25 MHz และ 2-FF Reset Synchronizer}

\subsubsection{การซิงโครไนซ์ปุ่มกด (2-FF Button Synchronizer)}
\todo{การซิงโครไนซ์ BTNU/D/L/R/C เข้าสู่ Clock Domain 25 MHz และการตรวจจับขอบสัญญาณ}

\subsubsection{การสร้างสัญญาณภาพ VGA (\texttt{vga\_sync.v} และ Output Register)}
\todo{ไทม์มิง VGA 640x480 @ 60 Hz, tick (y=480), tick2 (y=500) และ Output Register}

\subsubsection{การแสดงผลบน 7-Segment 8 หลัก (\texttt{seg7x8.v}) และ LEDs}
\todo{การแบ่งแสดงผลคะแนน (4 หลักขวา), เวลาร่าง 3 (หลัก 4-5), ร่างปัจจุบัน (หลักซ้ายสุด) และ LED[2:0] / LED[15:6]}

\subsection{ระบบตรรกะเกม (Game Logic System — \texttt{slime\_logic.v})}

\subsubsection{การสร้างตัวเลขสุ่มด้วย 16-bit LFSR}
\todo{อธิบาย Linear Feedback Shift Register สำหรับสุ่มการเคลื่อนที่ของผีและตำแหน่งเกิดมงกุฎ}

\subsubsection{การเดินบนตารางและการตรวจกำแพง (Tile Alignment \& \texttt{maze\_rom})}
\todo{พิกัดพิกเซล (px, py), เงื่อนไข aligned ที่ทวีคูณของ 32 และ next\_tx/next\_ty}

\subsubsection{ระบบคะแนนแบบ BCD (\texttt{bcd\_add10}) และตัวจับเวลา King Mode}
\todo{การบวกคะแนนแบบ BCD 4 หลัก และการนับถอยหลัง 10 วินาทีในร่าง 3}

\subsubsection{การตรวจจับการชน (Collision Detection)}
\todo{การตรวจสอบระยะห่างระหว่างสไลม์กับผี, เม็ดคะแนน และมงกุฎ}

\subsection{ระบบกราฟิกและการจัดการหน่วยความจำ (Graphic System)}

\subsubsection{แผนที่เขาวงกตและฟอนต์ (\texttt{slime\_roms.v})}
\todo{โครงสร้าง maze\_rom (20x15 ช่อง) และ font\_rom (ฟอนต์ 5x7)}

\subsubsection{ระบบสไปรต์และพาเลตสี (\texttt{gen\_art.py} และ \texttt{slime\_art.v})}
\todo{กระบวนการแปลงไฟล์ PNG จากโฟลเดอร์ sprite/ ด้วย gen\_art.py เป็น art\_palette (5-bit) และ ghost\_palette (4-bit) พร้อมตารางรวมสไปรต์จาก images/sprite/}

\subsubsection{การวาดภาพแบบ Priority Multiplexer (\texttt{slime\_draw.v})}
\todo{ลำดับชั้นการซ้อนทับพิกเซล: ข้อความ > สไลม์ > ผี > มงกุฎ > เม็ด > ไอคอน HUD > กำแพง > พื้น}

\subsection{ระบบจำลองการทำงานบนคอมพิวเตอร์ (Simulation System — Verilator + SDL2)}

\subsubsection{โหมดจำลอง 60 FPS (\texttt{fast\_top.v}, \texttt{sim\_fast.cpp})}
\todo{การจำลอง logic และ draw ที่ 60 เฟรมต่อวินาทีด้วย Verilator + SDL2}

\subsubsection{โหมดจำลองสัญญาณ VGA (\texttt{sim\_top.v}, \texttt{sim\_main.cpp})}
\todo{การจำลองสัญญาณ vga\_sync ทีละพิกเซลเสมือนฮาร์ดแวร์จริง}

% ============================================================================
\clearpage
\mainsection{ผลลัพธ์การทำงาน}
% ============================================================================

จากการออกแบบ พัฒนา และทดสอบระบบเกม Slime Maze ทั้งในตัวจำลอง Verilator + SDL2 และการรันจริงบนบอร์ด Digilent Nexys A7-100T เชื่อมต่อออกทางจอมอนิเตอร์ VGA ได้ผลลัพธ์การทำงานครบถ้วนตามคุณสมบัติที่กำหนด ดังแสดงในตารางที่ \ref{tab:results}

% [INSERT PICTURE: result_title.png, result_form1.png, result_form2_crown.png, result_king_mode.png, result_eat_ghost.png, result_pause.png, result_dead.png, result_win_over.png, result_board_hw.png]
\begin{longtable}{@{}p{0.48\textwidth} >{\raggedright\arraybackslash}p{0.46\textwidth}@{}}
    \caption{ผลลัพธ์การทำงานและการทดสอบระบบ}
    \label{tab:results} \\
    \toprule
    \textbf{ภาพผลลัพธ์การทำงาน} & \textbf{คำอธิบาย} \\
    \midrule
    \endfirsthead

    \multicolumn{2}{c}{\textit{(ต่อจากหน้าก่อน)}} \\
    \toprule
    \textbf{ภาพผลลัพธ์การทำงาน} & \textbf{คำอธิบาย} \\
    \midrule
    \endhead

    \midrule
    \multicolumn{2}{r}{\textit{(มีต่อหน้าถัดไป)}} \\
    \bottomrule
    \endfoot

    \bottomrule
    \endlastfoot

    % --- แถวที่ 1: Title ---
    \begin{minipage}[t]{\linewidth}
        \vspace{0pt}\centering
        \smartimage{0.96\linewidth}{result_title.png}{3.2cm}{ภาพหน้าจอ VGA หรือจากตัวจำลองแสดงกรอบคำว่า “SLIME MAZE” สีเหลืองทอง และ “PRESS BTNC” กะพริบอยู่กลางเขาวงกต}
        \tablefigcap{สาธิตการทำงานหน้าจอเริ่มต้น (\texttt{S\_TITLE})}
    \end{minipage}
    &
    \begin{minipage}[t]{\linewidth}
        \vspace{0pt}\raggedright
        เมื่อจ่ายไฟและโปรแกรม Bitstream หรือกดปุ่ม \texttt{CPU\_RESETN} เกมจะเข้าสู่สถานะ \texttt{S\_TITLE} หน้าจอแสดงแผนที่เขาวงกต $20\times15$, เม็ดคะแนน, สไลม์ร่างที่ 1, ผีทั้ง 4 ตัว และกล่องข้อความ “SLIME MAZE” พร้อมคำสั่ง “PRESS BTNC” กะพริบเพื่อรอรับการกดปุ่มกลางเริ่มเกม
    \end{minipage} \\
    \midrule

    % --- แถวที่ 2: Form 1 ---
    \begin{minipage}[t]{\linewidth}
        \vspace{0pt}\centering
        \smartimage{0.96\linewidth}{result_form1.png}{3.2cm}{ภาพขณะสไลม์ร่างที่ 1 กำลังเดินเก็บเม็ดในเขาวงกต มีเส้นทางที่เม็ดถูกกินไปแล้วบางส่วน และไอคอนชีวิต 3 ดวงมุมซ้ายบน}
        \tablefigcap{สาธิตการทำงานขณะเล่นในร่างที่ 1 (Form 1)}
    \end{minipage}
    &
    \begin{minipage}[t]{\linewidth}
        \vspace{0pt}\raggedright
        เมื่อกดปุ่ม \texttt{BTNC} เกมจะเข้าสู่สถานะ \texttt{S\_PLAY} ผู้เล่นควบคุมสไลม์ร่างที่ 1 เดินเก็บเม็ดคะแนน (+10 คะแนนต่อเม็ด) ด้วยปุ่ม \texttt{BTNU/D/L/R} โดยมีระบบกดเลี้ยวล่วงหน้า และหลบหลีกผีทั้ง 4 ตัว มุมซ้ายบนแสดงไอคอนชีวิต 3 ดวงเต็ม และ 7-Segment แสดงเลขร่าง “1” พร้อมคะแนนปัจจุบัน
    \end{minipage} \\
    \midrule

    % --- แถวที่ 3: Form 2 + Crown ---
    \begin{minipage}[t]{\linewidth}
        \vspace{0pt}\centering
        \smartimage{0.96\linewidth}{result_form2_crown.png}{3.2cm}{ภาพขณะสไลม์เปลี่ยนเป็นร่างที่ 2 (Big Slime) และเห็นไอเทมมงกุฎปรากฏอยู่ที่มุมหนึ่งของเขาวงกต}
        \tablefigcap{สาธิตการพัฒนาร่างเป็นร่างที่ 2 และการเกิดมงกุฎ}
    \end{minipage}
    &
    \begin{minipage}[t]{\linewidth}
        \vspace{0pt}\raggedright
        เมื่อเก็บเม็ดคะแนนครบ 30 เม็ด (\texttt{N\_EVOLVE = 30}) สไลม์จะวิวัฒนาการเป็นร่างที่ 2 (Big Slime) ที่มีขนาดและรูปลักษณ์เปลี่ยนไป พร้อมกับการปรากฏของไอเทมมงกุฎแอนิเมชันที่มุมใดมุมหนึ่งของแผนที่ เพื่อรอให้ผู้เล่นเดินไปเก็บและเข้าสู่โหมดราชา
    \end{minipage} \\
    \midrule

    % --- แถวที่ 4: King Mode ---
    \begin{minipage}[t]{\linewidth}
        \vspace{0pt}\centering
        \smartimage{0.96\linewidth}{result_king_mode.png}{3.2cm}{ภาพขณะสไลม์สวมมงกุฎ (ร่าง 3) ผีทุกตัวกลายเป็นสีฟ้า (\texttt{ghost\_scared}) และมุมขวาบนแสดงไอคอนรูปคิงโหมด (\texttt{kingmode})}
        \tablefigcap{สาธิตโหมดราชา (Form 3) และผีในสถานะหวาดกลัว}
    \end{minipage}
    &
    \begin{minipage}[t]{\linewidth}
        \vspace{0pt}\raggedright
        เมื่อเก็บมงกุฎได้ สไลม์จะเปลี่ยนเป็นร่างที่ 3 (King Slime) นาน 10 วินาที เดินเร็วขึ้น 2 เท่า ผีทุกตัวเปลี่ยนเป็นสไปรต์สีฟ้า (\texttt{ghost\_scared}) และวิ่งหนีออกจากสไลม์ ส่วนมุมขวาบนของจอภาพ (ช่อง \texttt{tx=19, ty=0}) จะแสดงไอคอน King Mode
    \end{minipage} \\
    \midrule

    % --- แถวที่ 5: Eating Ghosts & Blinking ---
    \begin{minipage}[t]{\linewidth}
        \vspace{0pt}\centering
        \smartimage{0.96\linewidth}{result_eat_ghost.png}{3.2cm}{ภาพขณะผีบางตัวถูกกินหายไปแล้ว (เหลือผีน้อยกว่า 4 ตัว) หรือจังหวะที่ผีกะพริบสลับสีเดิมในช่วง 2 วินาทีสุดท้ายของคิงโหมด}
        \tablefigcap{สาธิตการกินผีและการกะพริบเตือนท้ายโหมดราชา}
    \end{minipage}
    &
    \begin{minipage}[t]{\linewidth}
        \vspace{0pt}\raggedright
        ขณะอยู่ในโหมดราชา เมื่อสไลม์วิ่งชนผี ผีตัวนั้นจะหายไปจากสนามทันที (\texttt{galive[k] = 0}) และได้รับโบนัส +200 คะแนนต่อตัว และเมื่อเหลือเวลา $\le 2$ วินาทีสุดท้าย ตัวสไลม์ ผี และไอคอนคิงโหมดจะกะพริบเตือนก่อนกลับสู่ร่างที่ 2 และมงกุฎจะเกิดใหม่ในอีก 3 วินาที
    \end{minipage} \\
    \midrule

    % --- แถวที่ 6: Pause ---
    \begin{minipage}[t]{\linewidth}
        \vspace{0pt}\centering
        \smartimage{0.96\linewidth}{result_pause.png}{3.2cm}{ภาพหน้าจอกลางเขาวงกตแสดงข้อความ “PAUSE” สีเหลืองทอง โดยสไลม์และผีหยุดนิ่งอยู่กับที่}
        \tablefigcap{สาธิตการทำงานขณะหยุดเกมชั่วคราว (\texttt{S\_PAUSE})}
    \end{minipage}
    &
    \begin{minipage}[t]{\linewidth}
        \vspace{0pt}\raggedright
        เมื่อกดปุ่ม \texttt{BTNC} ระหว่างเล่นในสถานะ \texttt{S\_PLAY} เกมจะเข้าสู่สถานะ \texttt{S\_PAUSE} โดยหยุดการเคลื่อนที่ของสไลม์ ผี และตัวจับเวลาคิงโหมดทั้งหมด พร้อมแสดงกล่องข้อความ “PAUSE” และเมื่อกด \texttt{BTNC} อีกครั้งจะกลับมาเล่นต่อจากตำแหน่งเดิมทันที
    \end{minipage} \\
    \midrule

    % --- แถวที่ 7: Dead ---
    \begin{minipage}[t]{\linewidth}
        \vspace{0pt}\centering
        \smartimage{0.96\linewidth}{result_dead.png}{3.2cm}{ภาพขณะสไลม์แสดงท่าตาย (\texttt{pose=3}) และไอคอนชีวิตมุมซ้ายบนลดลงจาก 3 ดวงเป็น 2 ดวง (มี \texttt{life\_died} ปรากฏ 1 ดวง)}
        \tablefigcap{สาธิตการทำงานเมื่อสไลม์ชนผีและเสียชีวิต (\texttt{S\_DEAD})}
    \end{minipage}
    &
    \begin{minipage}[t]{\linewidth}
        \vspace{0pt}\raggedright
        หากสไลม์ในร่างที่ 1 หรือ 2 ชนเข้ากับผี ระบบจะเข้าสู่สถานะ \texttt{S\_DEAD} หน่วงเกม 1.5 วินาที (90 เฟรม) แสดงสไปรต์สไลม์ท่าตาย (\texttt{slime*\_died}) และไอคอนชีวิตมุมซ้ายบนจะเปลี่ยนเป็นสีเทา (\texttt{life\_died}) ตามจำนวนชีวิตที่เสียไป
    \end{minipage} \\
    \midrule

    % --- แถวที่ 8: Win & Over ---
    \begin{minipage}[t]{\linewidth}
        \vspace{0pt}\centering
        \smartimage{0.96\linewidth}{result_win_over.png}{3.2cm}{ภาพคู่เปรียบเทียบแสดงหน้าจอ “YOU WIN” สีเขียว และหน้าจอ “GAME OVER” สีแดง พร้อมข้อความ “PRESS BTNC”}
        \tablefigcap{สาธิตหน้าจอเมื่อชนะเกม (YOU WIN) และจบเกม (GAME OVER)}
    \end{minipage}
    &
    \begin{minipage}[t]{\linewidth}
        \vspace{0pt}\raggedright
        เมื่อกินผีได้ครบทั้ง 4 ตัว เกมจะเข้าสู่สถานะ \texttt{S\_WIN} แสดงข้อความสีเขียว “YOU WIN” แต่หากพลังชีวิตลดลงจนเหลือ 0 เกมจะเข้าสู่สถานะ \texttt{S\_OVER} แสดงข้อความสีแดง “GAME OVER” ทั้งสองสถานะสามารถกดปุ่ม \texttt{BTNC} เพื่อกลับไปยังหน้าจอ \texttt{S\_TITLE} และเริ่มรอบใหม่ได้
    \end{minipage} \\
    \midrule

    % --- แถวที่ 9: Hardware Board ---
    \begin{minipage}[t]{\linewidth}
        \vspace{0pt}\centering
        \smartimage{0.96\linewidth}{result_board_hw.png}{3.2cm}{ภาพถ่ายบอร์ด Nexys A7-100T จริงขณะเล่นเกม โดยเฉพาะจังหวะที่อยู่ในโหมดราชา แสดงตัวเลขบน 7-Segment และไฟ LED แถบเวลา LD15--LD6 รวมถึง LD2--LD0}
        \tablefigcap{ภาพการทำงานบนบอร์ดจริง แสดงผล 7-Segment และ LED}
    \end{minipage}
    &
    \begin{minipage}[t]{\linewidth}
        \vspace{0pt}\raggedright
        บนบอร์ด Digilent Nexys A7-100T จริง จอแสดงผล 7-Segment แสดงร่างปัจจุบันที่หลักซ้ายสุด, เวลานับถอยหลังโหมดราชาที่หลัก 5--4 และคะแนน BCD 4 หลักขวาสุด ขณะเดียวกัน \texttt{LED[2:0]} แสดงจำนวนชีวิตที่เหลือ และ \texttt{LED[15:6]} แสดงแถบเวลาโหมดราชาลดลงทีละดวงทุกวินาทีอย่างถูกต้อง
    \end{minipage} \\
\end{longtable}

% ============================================================================
\mainsection{ปัญหาอุปสรรคที่พบและแนวทางการแก้ไข}
% ============================================================================

\subsection{ปัญหาที่ 1 (เช่น การจัดการ Timing / การซิงโครไนซ์สัญญาณ)}
\todo{ระบุปัญหาที่พบ, สาเหตุ และแนวทางการแก้ไข}

\subsection{ปัญหาที่ 2 (เช่น ข้อจำกัดของ Palette หรือการควบคุม)}
\todo{ระบุปัญหาที่พบ, สาเหตุ และแนวทางการแก้ไข}

% ============================================================================
\mainsection{สรุป}
% ============================================================================
\todo{สรุปผลการดำเนินโครงงานและองค์ความรู้ที่ได้รับจากรายวิชา CPE222}

% ============================================================================
\clearpage
\mainsection{บรรณานุกรม}
% ============================================================================
\todo{รายการเอกสารอ้างอิงและแหล่งข้อมูลที่ใช้ในโครงงาน (เช่น Digilent Nexys A7-100T Reference Manual, VGA Timing, Verilator, SDL2)}

% ============================================================================
\clearpage
\mainsection{ภาคผนวก}
% ============================================================================
\todo{ลิงก์ซอร์สโค้ด Verilog (GitHub Repository) และวิดีโอสาธิตการทำงาน}


\end{document}
